
\section*{0.5. Base Change}
\phantomsection
\label{sga3:VIIB:0.5}
\addcontentsline{toc}{subsection}{\SGARefLink{sga3:VIIB:0.5}{0.5}. Base Change}

Let
\[
  \phi:k\longrightarrow\ell
\]
be a homomorphism of pseudocompact rings (\SGARefLink{sga3:VIIB:0.1.3}{0.1.3}). The construction of
\SGARefLink{sga3:VIIB:0.3}{0.3} admits the following generalization.\footnote{The editors expanded
the details of this paragraph.}

\medskip
\noindent\textbf{Definition 0.5.A.}
\phantomsection
\label{sga3:VIIB:0.5.A}
{\itshape For every \(M\in\operatorname{PC}(k)\) and
\(N\in\operatorname{PC}(\ell)\), let
\[
  M\widehat{\otimes}_kN
\]
be the separated completion of \(M\otimes_kN\) for the linear topology
defined by the kernels of the projections
\[
  M\otimes_kN
  \longrightarrow
  (M/M')\otimes_k(N/N'),
\]
where \(M'\), respectively \(N'\), ranges over the open submodules of
\(M\), respectively \(N\). Then \(M\widehat{\otimes}_kN\) is a
pseudocompact \(\ell\)-module. If \(m\in M\) and \(x\in N\), the
canonical image of \(m\otimes_kx\) in
\(M\widehat{\otimes}_kN\) is denoted by
\[
  m\widehat{\otimes}_kx.
\]}

This applies in particular to \(N=\ell\). In that case
\[
  M\widehat{\otimes}_k\ell
\]
is called the pseudocompact \(\ell\)-module obtained from \(M\) by the
base change \(k\to\ell\).

\medskip
\noindent\textbf{Remarks 0.5.B.}
\phantomsection
\label{sga3:VIIB:0.5.B}

\noindent\textit{\SGARefLink{sga3:vii:paragraph:5-1:equation:i}{(i)}}
In such a base change, \(\ell\) need not be a profinite
\(k\)-algebra. The typical example is that in which \(k\) is a field
and \(\ell=K\) is an arbitrary field extension of \(k\).

\smallskip
\noindent\textit{\SGARefLink{sga3:vii:paragraph:5-1:equation:ii}{(ii)}}
If, however, the \(k\)-module underlying \(N\) is pseudocompact---for
example, if \(\ell\) is a profinite \(k\)-algebra---then \SGARefLink{sga3:VIIB:0.4}{0.4} shows
that every open \(k\)-submodule of \(N\) contains an open
\(\ell\)-submodule. In that case \(M\widehat{\otimes}_kN\) agrees with
the completed tensor product of \SGARefLink{sga3:VIIB:0.3}{0.3} for the pseudocompact \(k\)-modules
\(M\) and \(N\), so the notation is unambiguous.

In every case, the \(k\)-module \(N|_k\) obtained by restriction of
scalars is a topological \(k\)-module: the map
\[
  k\times N\longrightarrow N,
  \qquad
  (t,n)\longmapsto\phi(t)n
\]
is continuous. We write
\(\operatorname{Hom}_c(M,N|_k)\) for the abelian group of continuous
\(k\)-linear maps from \(M\) to \(N|_k\).

\medskip
\noindent\textbf{Proposition 0.5.C.}
\phantomsection
\label{sga3:VIIB:0.5.C}
\textit{For every \(M\in\operatorname{PC}(k)\) and
\(N\in\operatorname{PC}(\ell)\), there is a functorial isomorphism}
\[
  \operatorname{Hom}_{\operatorname{PC}(\ell)}
    (M\widehat{\otimes}_k\ell,N)
  \simeq
  \operatorname{Hom}_c(M,N|_k).
\]

Indeed, let \(u:M\to N|_k\) be a continuous homomorphism. The map
\[
  u':M\times\ell\longrightarrow N,
  \qquad
  (m,\lambda)\longmapsto\lambda u(m)
\]
is continuous and ``bilinear,'' meaning \(k\)-linear in the first
factor and \(\ell\)-linear in the second. If \(N'\) is an open
\(\ell\)-submodule of \(N\), the argument of \SGARefLink{sga3:VIIB:0.3.1}{Lemma 0.3.1} gives an open
submodule \(M'\subset M\) and an open ideal
\(\ell'\subset\ell\) such that
\[
  u'(M'\times\ell)\subset N',
  \qquad
  u'(M\times\ell')\subset N'.
\]
Thus \(u\) induces a continuous \(\ell\)-linear map
\[
  U:M\widehat{\otimes}_k\ell\longrightarrow N
\]
satisfying
\[
  U(m\widehat{\otimes}_k\lambda)=\lambda u(m).
\]
Conversely, every morphism
\[
  f:M\widehat{\otimes}_k\ell\longrightarrow N
\]
gives the morphism
\[
  f':M\longrightarrow N|_k,
  \qquad
  m\longmapsto f(m\widehat{\otimes}_k1).
\]

As in \SGARefLink{sga3:VIIB:0.3.2}{0.3.2}, \SGARefLink{sga3:VIIB:0.3.5}{0.3.5}, and \SGARefLink{sga3:VIIB:0.3.1.2}{0.3.1.2}, this yields:

\medskip
\noindent\textbf{Corollary 0.5.D.}
\phantomsection
\label{sga3:VIIB:0.5.D}
{\itshape The functor
\[
  \operatorname{PC}(k)\longrightarrow\operatorname{PC}(\ell),
  \qquad
  M\longmapsto M\widehat{\otimes}_k\ell
\]
is right exact and commutes with filtered projective limits. Thus, if
\((M_i)\) is a filtered projective system in
\(\operatorname{PC}(k)\), there is a canonical isomorphism
\[
  (\varprojlim_iM_i)\widehat{\otimes}_k\ell
  \simeq
  \varprojlim_i(M_i\widehat{\otimes}_k\ell).
\]
Moreover, if \(M\) and \(N\) are pseudocompact \(k\)-modules, there is
a canonical isomorphism
\[
  (M\widehat{\otimes}_kN)\widehat{\otimes}_k\ell
  \simeq
  (M\widehat{\otimes}_k\ell)
    \widehat{\otimes}_{\ell}
  (N\widehat{\otimes}_k\ell).
\]}

\medskip
\noindent\textbf{Definition 0.5.E.}
\phantomsection
\label{sga3:VIIB:0.5.E}
{\itshape Finally, if \(A\) is a profinite \(k\)-algebra, there is a
unique profinite \(\ell\)-algebra structure on
\(A\widehat{\otimes}_k\ell\) such that
\[
  (a\widehat{\otimes}_k\lambda)
  (b\widehat{\otimes}_k\mu)
  =
  ab\widehat{\otimes}_k(\lambda\mu)
\]
for \(a,b\in A\) and \(\lambda,\mu\in\ell\). The profinite
\(\ell\)-algebra \(A\widehat{\otimes}_k\ell\) is said to be obtained
from \(A\) by extension of scalars, or base change, along
\(k\to\ell\).}

